
\documentclass[../main.tex]{subfiles}
\graphicspath{{\subfix{../images/}}}

\begin{document}
\part{Part4几何}

\section{init}
到一个定点距离为定值的点的轨迹是 圆

到两个定点距离相等的点的轨迹是 直线(垂直平分线)

到两个定点距离之比为定值($\ne 1$)的点的轨迹是圆，也就是阿波罗尼奥斯圆，简称阿氏圆

到两个定点距离之和为定值的点的轨迹是 椭圆

到两个定点距离之差为定值的点的轨迹是 双曲线

到定点与定直线距离相等的点的轨迹是 抛物线

\clearpage
\section{三角形}

\textbf{大边对大角、大角对大边}，是定理吗？其实这个命题出现在欧几里得《几何原本》第一卷Proposition 18(In any triangle the greater side subtends the greater angle)、19(In any triangle the greater angle is subtended by the greater  side)，
欧几里得先证明了大边对大角(取等腰)，然后用大边用大角的结论，反证了大角对大边，哎呀，用了反证法，欧几里得的最爱。

\textbf{Proposition 06:} 等角对等边

\textbf{Proposition 16:} 外角大于任一内对角

\textbf{Proposition 20:} 任意三角形中，任意两边之和大于第三边


\begin{minipage}{0.7\textwidth}
  \textbf{Proposition 21:} 若从三角形一边的两个端点作两条直线交于三角形内，则这样作出的两条直线之和小于三角形其余两边之和，但夹角更大。如右图：
  $BD+CD < AB+AC, \; \angle BDC> \angle BAC$.\\
  作辅助线$BD$延长线交$AC$与$E$，用极限的思想思考，$D$无限接近于$A$，这时候$DE$无限接近于$0$，$DE$就是突破口。\\
  $AB+AE>BD+DE; \; CE+DE>CD$得证。$\angle BDC>\angle BEC>\angle BAC$得证。
\end{minipage}
\begin{minipage}{0.3\textwidth}
  \centering
  \begin{tikzpicture}
    \tkzDefPoints{1/2.5/A, 0/0/B, 3/0/C, 1.1/1/D}
    \tkzInterLL(B,D)(A,C) \tkzGetPoint{E}

    \tkzDrawSegments(A,B A,C B,C D,C B,D)
    \tkzDrawSegments[dashed](D,E)
    \tkzLabelPoints[above](A)
    \tkzLabelPoints[left](B)
    \tkzLabelPoints[right](C)
    \tkzLabelPoints[above left](D)
    \tkzLabelPoints[above right](E)
  \end{tikzpicture}
\end{minipage}

\textbf{四心:}内心(内切圆)：角平分线; 外心(外接圆)：垂直平分线; 垂心： 垂线 ;重心： 中线。

\subsection{重要定理}
\begin{minipage}{0.6\textwidth}
  \textbf{余弦定理：}此定理可追溯至《几何原本》，第二卷命题12、13，书中分为钝角和锐角三角形。如右图：\\
  钝角：$BC^2=AC^2+AB^2+2AD\cdot AB$\\
  锐角：$BC^2=AC^2+AB^2-2AD\cdot AB$
\end{minipage}
\begin{minipage}{0.4\textwidth}
  \centering
  \begin{tikzpicture}[scale=1]
    \tkzDefPoints{.7/0/A, 2.4/0/B, 0/2/C, 0/0/D}
    \tkzDrawSegments(A,B A,C B,C)
    \tkzDrawSegments[dashed](A,D C,D)
    \tkzLabelPoints(A,B)
    \tkzLabelPoints[left](C,D)
  \end{tikzpicture}
  \begin{tikzpicture}[scale=1]
    \tkzDefPoints{1/0/A, 3/0/B, 1.7/2/C, 1.7/0/D}
    \tkzDrawSegments(A,B A,C B,C)
    \tkzDrawSegments[dashed](A,D C,D)
    \tkzLabelPoints(A,B,D)
    \tkzLabelPoints[left](C)
  \end{tikzpicture}
\end{minipage}

现代数学式中$AD$分别为：$AC\cos(\pi - \angle BAC)$, $AC\cos\angle BAC$，即为余弦定理形式。

任意$\triangle ABC$，$a,b,c$分别为$\angle A, \angle B, \angle C$的对边，则有：\\
$c^2=a^2+b^2-2ab\cos \angle C \phantom{xx} b^2=a^2+c^2-2ac\cos \angle B \phantom{xx} a^2=b^2+c^2-2bc\cos \angle A $

证明过程也很简单，两个直角三角形运用勾股定理即可。此定理主要解决非直角三角形的“解三角形”问题：1.$SAS$求第三边；2.$SSS$求任意角；3.判断三角形形状(钝角、锐角、直角)；4.工程测量

\textbf{正弦定理：}任意$\triangle ABC$，$a,b,c$分别为$\angle A, \angle B, \angle C$的对边，$R$为外接圆半径，则有：
$$\frac{a}{\sin \angle A} = \frac{b}{\sin \angle B} = \frac{c}{\sin \angle C}=2R $$
主要应用：1.一直两个角和任意一边，可以求出另外两边；2.一直两边和其中一边的对角，可以求出其他角和边(需判断解的数量); 3.求外接圆半径。

思考：正弦定理可以用来证明 大边对大角、大角对大边吗？

\begin{minipage}{0.7\textwidth}
  \textbf{斯图尔特定理(Stewart's Theorem)：}任意三角形为$ABC$，$D$为$BC$上一点，已知三边长和$D$的截取底边位置，求$AD$的长度。如果有$AB$列式子，会比较凌乱，不如设$AB=a, AC=b, BC=c, BD=x, DC=y, AD=t$，之后再代回：
\end{minipage}
\begin{minipage}{0.3\textwidth}
  \centering
  \begin{tikzpicture}[scale=.9]
    \tkzDefPoints{1.2/2/A, 0/0/B, 3/0/C, 1.4/0/D}
    \tkzDrawSegments(A,B A,C B,C A,D)
    \tkzLabelPoints[above](A)
    \tkzLabelPoints[left](B)
    \tkzLabelPoints[right](C)
    \tkzLabelPoints(D)
    \tkzLabelSegment[left, midway, red](A,B){a}
    \tkzLabelSegment[right, midway, red](A,C){b}
    \tkzLabelSegment[below, midway, red](B,D){x}
    \tkzLabelSegment[below, midway, red](D,C){y}
    \tkzLabelSegment[left, midway, red](A,D){t}
  \end{tikzpicture}
\end{minipage}
\[
  \left\{
    \begin{aligned}
      a^2&=x^2+t^2-2xt\cos\angle BDA\\
      b^2&=y^2+t^2-2yt\cos(180^{\circ}-\angle BDA)
    \end{aligned}
  \right.
  \Rightarrow
    t^2(x+y)+xy(x+y)=a^2y+b^2x
\]
当$x=y$时，也即$AD$是中线时，也就是中位线定理：$2t^2+2x^2=a^2+b^2$\\
当$x/y=a/b$时，也就是$AD$是角平分线时，此时是内角平分线长度：$t^2=ab-xy$。

如果$a=b, \; x=y$，也就是三角形为等腰时，$D$为底边中点时，等式就是勾股定理。

这个定理，没必要去记忆，记忆也很困难。但是需要知道推导过程，这样以来，中线定理，角平分线的长度很容易推导出来。

\textbf{射影定理:}又称欧几里得定理(Euclid's theorem)，Geometric Mean Theorem。在一个直角三角形里，一条直角边的平方等于三角形的斜边乘以该直角边在斜边上的正投影。此定理出现在《几何原本》第一卷毕氏定理的证明过程中。

\begin{minipage}{.4\textwidth}
  \centering
  \begin{tikzpicture}[scale=1,rotate=-123]
    \tkzDefPoints{0/0/C, 0/3/B, 2/0/A}
    \tkzDefPointsBy[projection=onto A--B](C){D}
    \tkzDrawPolygons(A,B,C)
    \tkzDrawSegments(C,D)
    \tkzMarkRightAngles[scale=0.8, color=teal](C,D,A A,C,B)
    \tkzLabelPoints(D)
    \tkzLabelPoints[left](A,C)
    \tkzLabelPoints[right](B)
  \end{tikzpicture}
\end{minipage}
\begin{minipage}{.6\textwidth}
  $AC^2=AD\cdot AB$\\
  $BC^2=BD\cdot AB$\\
  $CD^2=AD\cdot BD${
    \footnotesize
    (几何平均定理，几何原本第六卷第8个命题)
  }
\end{minipage}

$AD,BD$为$AC,BC$在底边$AB$上的正投影，因此得名。三个相似的直角三角形可证。\\前两个等式相加便是勾股定理。

\subsection{全等}
SAS中的S表示Side, A为Angle

几何原本第一卷关于三角形全等集中在命题4、8、26。命题4即是SAS，用“移动重合” 的方法证明。命题8为SSS。命题26为ASA、AAS。

\begin{minipage}{0.7\textwidth}
  全等除了AAA(相似)和ASS(屁股)(如右图)，其余都是全等。

  SSS, SAS, AAS(ASA,SAA)

  对于直角三角形全等判定，因为多了一个特定的直角，由勾股定理，已知两边也就确定了第三边，SSA也可以，也就是HL。
\end{minipage}
\begin{minipage}{0.3\textwidth}
  \centering
  \begin{tikzpicture}
    \tkzDefPoints{2/2/A, 0/0/B, 3/0/C, 1/0/D}
    \tkzDrawSegments[thick, red](A,B)
    \tkzDrawSegments(B,C)
    \tkzDrawSegments[dashed, red, thick](A,C A,D)
    \tkzMarkAngles[scale=0.6](C,B,A)
    \tkzFillAngles[scale=0.6, fill=red!20,opacity=.5](C,B,A)
    \tkzMarkSegments[mark=||, size=3pt](A,C A,D)
  \end{tikzpicture}
\end{minipage}

\subsection{相似}
几何原本第六卷，相似三角形的定义：如果两个三角形对应的角相等，且对应边的比相等，则这两个三角形相似。

命题2：平行于三角形底边的直线截得的三角形两边的线段成比例。\\
命题3：三角形角平分线截得的底边线段之比等于两侧边之比。\\
命题4：(AAA)如果两个三角形各角相等，那么等角所对应的边成比例。\\
命题5：如果两个三角形的边成比例，那么各边所对的角相等。\\
命题6：(SAS)如果两个三角形有一个角相等，且夹该等角的边成比例，那么这两个三角形各角对应相等。\\
命题7：(ASS带限定条件)如果两个三角形有一个角相等，且夹另外一个角的边成比例，其余的角要么都小于要么都不小于一直角，则这两个三角形是等角的，成比例的边所夹的角相等。\\
命题8：直角三角形

命题21：(相似的传递性)与同一直线形相似的图形也彼此相似

\textbf{判定定理}有三个：1.$AAA$; 2.$SAS$; 3.$SSS$。

很明显，相似的三角形，周长比等于相似比；面积比等于相似比的平方(海伦公式)。

$\triangle ABC \sim \triangle A'B'C'$, 字母顺序是对应的：$AB$一定对应$A'B'$，$\angle ABC$对应$\angle A'B'C'$。相似又分为顺相似和逆相似，逆相似一般需要镜像、翻转和原图像重合，顺相似只需要平移、旋转，这两个概念一般不做要求。

相似的比例有两种写法：$\displaystyle \frac{AB}{AC}=\frac{A'B'}{A'C'},  \text{or}\;\; \frac{AB}{A'B'}=\frac{AC}{A'C'}=\frac{BC}{B'C'}(=k)$，注意微妙的不同，前者可以写为一个$ AB:AC:BC=DE:DF:EF$。几何原本里，用的是前者。

\begin{minipage}{0.25\textwidth}
  \centering
  \begin{tikzpicture}[scale=0.8]
    \tkzDefPoints{2/2/A, 0/0/B, 3/0/C, 2.4/0/D}
    \tkzDrawSegments(A,B A,C B,C)
    \tkzLabelPoints[left](B,A)
    \tkzLabelPoints[right](C)
  \end{tikzpicture}
\end{minipage}
\begin{minipage}{0.25\textwidth}
  \centering
  \begin{tikzpicture}[scale=0.6]
    \tkzDefPoints{2/2/A', 0/0/B', 3/0/C'}
    \tkzDrawPolygons(A',B',C')
    \tkzLabelPoints[left](A',B')
    \tkzLabelPoints[right](C')
  \end{tikzpicture}
\end{minipage}
\begin{minipage}{0.25\textwidth}
  \centering
  \begin{tikzpicture}[scale=0.8]
    \tkzDefPoints{2/2/A, 0/0/B, 3/0/C, 2.4/0/D}
    \tkzDrawSegments(A,B A,C B,C)
    \tkzDrawSegments[dashed](A,D)

    \tkzMarkAngles[size=0.4, mark=||, mksize=2pt](B,A,D)
    \tkzFillAngles[size=0.4, fill=red!20,opacity=.5](B,A,D)
    \tkzMarkAngles[size=0.4, mark=||, mksize=2pt](A,C,B)
    \tkzFillAngles[size=0.4, fill=red!20,opacity=.5](A,C,B)

    \tkzLabelPoints[left](B,A)
    \tkzLabelPoints[right](C)
    \tkzLabelPoints(D)
  \end{tikzpicture}
\end{minipage}
\begin{minipage}{0.25\textwidth}
  \centering
  \begin{tikzpicture}[scale=1]
  \tkzDefPoints{2/1/A, 0/0/B, 3/0/C, 1.8/0/D}
    \tkzDrawSegments(A,B A,C B,C)
    \tkzDrawSegments[dashed](A,D)

    \tkzMarkAngles[size=0.4, mark=||, mksize=2pt](B,A,D)
    \tkzFillAngles[size=0.4, fill=red!20,opacity=.5](B,A,D)
    \tkzMarkAngles[size=0.4, mark=||, mksize=2pt](A,C,B)
    \tkzFillAngles[size=0.4, fill=red!20,opacity=.5](A,C,B)

    \tkzLabelPoints[above](A)
    \tkzLabelPoints[left](B)
    \tkzLabelPoints[right](C)
    \tkzLabelPoints(D)
  \end{tikzpicture}
\end{minipage}

有一个比较特别的相似，就是三角形内部的相似，如上右两图：
$\triangle ABD \sim \triangle CBA$, 共用角$\angle A$，如何确定顺序？那么可以从共同的角开始，$\angle ABD$ 要么对应$\angle ABC$要么是$\angle CBA$，一般$AB$不是同一个位置，典型的逆相似，需要翻折，或者再找另一个角，所以……

\subsection{角平分线、角平分线定理、内心}
\begin{minipage}{0.5\textwidth}
  \centering
  \begin{tikzpicture}[rotate=-20]
    \tkzDefPoints{0/0/B, 2/0/C, 1.7/1.7/A}

    \tkzDefLine[bisector](B,A,C)  \tkzGetPoint{a}
    \tkzInterLL(B,C)(A,a) \tkzGetPoint{D}

    \tkzDefLine[parallel=through C](A,D) \tkzGetPoint{c'}
    \tkzInterLL(C,c')(B,A) \tkzGetPoint{C'}

    \tkzMarkAngles[scale=.8](B,A,D A,C',C D,A,C C',C,A)
    \tkzLabelAngles[pos=1, red, scale=0.8](B,A,D A,C',C){$\alpha$}
    \tkzLabelAngles[pos=1, red, scale=0.8](D,A,C C',C,A){$\beta$}

    \tkzDrawSegments(B,C A,C B,A)
    \tkzDrawSegments[dashed](A,D C,C' A,C')
    \tkzLabelPoints(B,C,D)
    \tkzLabelPoints[above left](A, C')
  \end{tikzpicture}

  内角平分线定理及其逆定理\\
  $\angle \alpha=\angle \beta \Leftrightarrow \frac{AB}{AC} = \frac{BD}{DC}$
\end{minipage}
\begin{minipage}{0.5\textwidth}
  \centering
  \begin{tikzpicture}
    \tkzDefPoints{0/0/B, 2/0/C, 2.8/2/A}

    \tkzDefLine[bisector out](B,A,C)  \tkzGetPoint{a'}
    \tkzInterLL(B,C)(A,a') \tkzGetPoint{E}

    \tkzDefPointOnLine[pos=1.5](B,A)  \tkzGetPoint{A'}

    \tkzDefLine[parallel=through C](A,E) \tkzGetPoint{c'}
    \tkzInterLL(C,c')(B,A) \tkzGetPoint{C'}

    \tkzMarkAngles[scale=0.6](C,A,E E,A,A' A,C,C' C,C',A)
    \tkzLabelAngles[pos=0.7, red, scale=0.8](C,A,E A,C,C'){$\gamma$}
    \tkzLabelAngles[pos=0.7, red, scale=0.8](E,A,A' C,C',A){$\delta$}

    \tkzDrawSegments(B,C A,C B,A)
    \tkzDrawSegments[dashed](A,E C,E A,A' C,C')
    \tkzLabelPoints(B,C,E)
    \tkzLabelPoints[above left](A, C')
  \end{tikzpicture}

  外角平分线定理及其逆定理\\
  $\angle \alpha=\angle \beta \Leftrightarrow \frac{AB}{AC} = \frac{BD}{DC}$
\end{minipage}

上面两图是从$C$作内角、外角平分线构造。内角平分线定理，也可以用面积相等来证明。\\
$S_{\triangle ADB}:S_{\triangle ADC}=\frac{1}{2}AB\cdot AD\sin \alpha:\frac{1}{2}AD\cdot AC\sin \beta=\frac{1}{2}BD\cdot h:\frac{1}{2}CD\cdot h$

\textbf{内角平分线的长度：}$AD^2=AB\cdot AC-BD\cdot DC$，$AB=AC$时，也就是勾股定理 \\
\textbf{外角平分线的长度：}$AE^2=BE\cdot CE-AB\cdot AC$

\textbf{阿波罗尼奥斯圆(Apollonius Circle 阿氏圆)}，角平分线性质有广泛的应用，其中一个关系相当紧密的应用就是阿氏圆，也就是\textbf{平面上到两定点距离之比为定值(不等于$1$)的点的轨迹}是一个圆，它的代数证明很简单，几何证明需要借助角平分线定理。圆的内分点和外分点为阿氏圆的直径。下图：

\begin{minipage}{1.0\textwidth}
  \centering
  \begin{tikzpicture}
    \tkzDefPoints{0/0/B, 2/0/C, 2.8/2/A}
    \tkzDefLine[bisector out](B,A,C)  \tkzGetPoint{a'}
    \tkzInterLL(B,C)(A,a') \tkzGetPoint{E}
    \tkzDefPointOnLine[pos=1.2](B,A)  \tkzGetPoint{A'}

    \tkzMarkAngles[scale=0.6, dashed](C,A,E E,A,A')
    \tkzLabelAngles[pos=0.7, red, scale=0.8](C,A,E E,A,A'){$\gamma$}
    \tkzDefLine[bisector](B,A,C)  \tkzGetPoint{a}
    \tkzInterLL(B,C)(A,a) \tkzGetPoint{D}
    \tkzMarkAngles[scale=0.8, dashed](B,A,D D,A,C)
    \tkzLabelAngles[pos=1, red, scale=0.8](B,A,D D,A,C){$\alpha$}
    \tkzDefCircle[circum](A,D,E) \tkzGetPoint{O}
    \tkzDrawCircle[dashed, thick](O,A)

    \tkzDrawSegments(B,C A,C B,A)
    \tkzDrawSegments[dashed](A,E C,E A,A' A,D)
    \tkzDrawPoints(D,E,O)
    \tkzLabelPoints(B,C,D,E,O)
    \tkzLabelPoints[above left](A)

  \end{tikzpicture}
\end{minipage}

例：同用上图，满足条件$BC=3, \; AB=2AC$的三角形的面积的最大值。\\
因为底边是个定值，面积最大只需要高最大即可，其余两边长的比例是定值，所以定点$A$在一个阿氏圆上，只需要求出半径即可。
$CD=1, \;BD=2$这个是内角平分线，再用一次外角平分线，$BE:CE=2:1 \Rightarrow (2r+BD):(2r-CD)=2:1$，可算出$r$。此题如果用代数方法，比如余弦定理或者海伦公式会很麻烦。

总体来说，阿氏圆有点抽象。已知底边长和两侧边之比，那么底边上的内分点$D$和外分点$E$就可以确定，那么这个点的轨迹连接$D,E$都可以形成直角(如上图，圆上任一点，连接$AD,AE$)，很明显$DE$是一个圆的直径。隐含有一对相似三角形: $S_{\triangle COA} \sim S_{\triangle AOB} \Leftrightarrow CO:AO=CA:AB=OA:OB \Leftrightarrow CO:r=CA:AB=r:OB$.

初中几何中的阿氏圆问题，一般是圆外两点到圆上的距离之和，如：$PA+\frac{1}{2}PB$，那么把系数小于$1$的定点$B$转化为圆内，这个定点$B'$在$A$到圆心的线段上。感觉用阿氏圆比较抽象，用相似三角形较为容易理解，但是，系数已知，圆已知，圆外的定点已知，所以题目一般是巧合的题目($r:OB$，比例已知)，是往阿氏圆上凑的题目。

\begin{minipage}{0.5\textwidth}
  \textbf{内心:} \\inscribed内切的、雕刻的. \\ Name was inscribed on the trophy
\end{minipage}
\begin{minipage}{0.5\textwidth}
  \begin{center}
    \begin{tikzpicture}
      \tkzDefPoints{1.2/2.1/A, 0/0/B, 3/0/C}

      \tkzDefCircle[in](A,B,C)   \tkzGetPoints{O}{x}
      \tkzDrawCircles(O,x)
      \tkzDefPointsBy[projection=onto A--B](O){D}
      \tkzDefPointsBy[projection=onto A--C](O){E}
      \tkzDefPointsBy[projection=onto B--C](O){F}
      \tkzMarkRightAngle[scale=0.6](O,D,A)
      \tkzMarkRightAngle[scale=0.6](O,E,A)
      \tkzMarkRightAngle[scale=0.6](O,F,B)
      \tkzMarkSegments[mark=|, size=2pt](O,D O,E O,F)

      \tkzDrawSegments(A,B A,C B,C O,A O,B O,C)
      \tkzDrawSegments[dashed](O,D O,E O,F)
      \tkzDrawPoints(O)
      \tkzLabelPoints[above](A)
      \tkzLabelPoints[left](B)
      \tkzLabelPoints[right](C,O)
    \end{tikzpicture}
  \end{center}
\end{minipage}

\subsection{中线、重心定理、中位线定理}
\begin{minipage}{0.33\textwidth}
  \centering
  \begin{tikzpicture}
    \tkzDefPoints{0/0/B, 3/0/C, 1.2/2.8/A}
    \tkzDefPointOnLine[pos=0.5](A,B)  \tkzGetPoint{F}
    \tkzDefPointOnLine[pos=0.5](B,C)  \tkzGetPoint{D}
    \tkzDefPointOnLine[pos=0.5](A,C)  \tkzGetPoint{E}
    \tkzInterLL(A,D)(B,E) \tkzGetPoint{G}

    \tkzDrawSegments(B,C A,C B,A G,A G,B G,C)
    \tkzDrawPoints[red,size=4](G)
    \tkzLabelPoints[above left](A)
    \tkzLabelPoints[left](B)
    \tkzLabelPoints[right](C,G)
  \end{tikzpicture}
\end{minipage}
\begin{minipage}{0.33\textwidth}
  \centering
  \begin{tikzpicture}
    \tkzDefPoints{0/0/B, 3/0/C, 1.2/2.8/A}
    \tkzDefPointOnLine[pos=0.5](A,B)  \tkzGetPoint{F}
    \tkzDefPointOnLine[pos=0.5](B,C)  \tkzGetPoint{D}
    \tkzDefPointOnLine[pos=0.5](A,C)  \tkzGetPoint{E}
    \tkzInterLL(A,D)(B,E) \tkzGetPoint{G}

    \tkzDrawSegments(B,C A,C B,A C,F B,E A,D)
    \tkzDrawPoints[red,size=4](G)
    \tkzLabelPoints[above left](A)
    \tkzLabelPoints[left](B)
    \tkzLabelPoints[right](C,G)
    \tkzLabelPoints(D)
  \end{tikzpicture}
\end{minipage}
\begin{minipage}{0.33\textwidth}
  \centering
  \begin{tikzpicture}
    \tkzDefPoints{0/0/B, 3/0/C, 1/2.8/A}
    \tkzDefPointOnLine[pos=0.5](A,B)  \tkzGetPoint{F}
    \tkzDefPointOnLine[pos=0.5](A,C)  \tkzGetPoint{E}
    \tkzDefPointOnLine[pos=0.5](B,C)  \tkzGetPoint{D'}
    \tkzInterLL(C,F)(B,E) \tkzGetPoint{G}

    \tkzDrawSegments(B,C A,C B,A C,F B,E)
    \tkzDefPointOnLine[pos=2](A,G)  \tkzGetPoint{H}

    \tkzDrawSegments[dashed](B,H H,C)
    \tkzDrawSegments[thick, dashed, red](A,G)
    \tkzDrawSegments[thick, dashed, blue](G,H)
    \tkzMarkSegments[mark=||](A,G G,H)
    \tkzDrawPoints(G)
    \tkzLabelPoints[above left](A,G)
    \tkzLabelPoints[left](B,F)
    \tkzLabelPoints[right](C,E)
    \tkzLabelPoints(H,D')
  \end{tikzpicture}
\end{minipage}

\textbf{重心定理：}三条中线相交于一点，此点到顶点的距离是它到对边中点距离的2倍($AG=2GD$)

\textbf{中位线定理：} $AB^2+AC^2=2(AD^2+BD^2)$，由两个余弦定理相加很容易得出

\textbf{面积特征：} $S_{\triangle AGB} = S_{\triangle AGC} = S_{\triangle BGC}=\frac{1}{3}S_{\triangle ABC}$

\textbf{坐标表示：}Let $A(x_1,y_1), \;B(x_2, y_2), \; C(x_3, y_3)$, then $\displaystyle G(\frac{x_1+x_2+x_3}{3}, \frac{y_1+y_2+y_3}{3})$,
Proof\footnote{$\vv{AG}=2\vv{GD}=\vv{GC}+\vv{GB} \Rightarrow (x-x_1, y-y_1)=(x_2-x+x_3-x, y_2-y+y_3-y)$}

\textbf{向量特征：}$\vv{OG} = \displaystyle \frac{1}{3}(\vv{OA}+\vv{OB}+\vv{OC})$，$O$是原点

\textbf{向量特征}：$\vv{GA} + \vv{GB} + \vv{GC}=0$

\textbf{物理意义：}对于质量均匀的薄片，重心为三中线的交点，重心是三角形的物理平衡点(质心)

\textbf{性质：}重心到三个顶点的距离的平方和最小。直接用距离公式平方相加。

\subsection{垂心}
\begin{minipage}{0.33\textwidth}
  \centering
  \begin{tikzpicture}[scale=1]
    \tkzDefPoints{1/2.7/A, 0/0/B, 3.6/0/C}
    \tkzDefPointsBy[projection=onto A--B](C){D}
    \tkzDefPointsBy[projection=onto A--C](B){E}
    \tkzDefPointsBy[projection=onto B--C](A){F}
    \tkzInterLL(A,F)(B,E) \tkzGetPoint{H}

    \tkzMarkRightAngles[scale=0.8](C,D,B B,E,C C,F,A)
    \tkzDrawSegments(A,B A,C B,C B,E C,D A,F)
    \tkzDrawPoints(H)
    \tkzLabelPoints[above](A)
    \tkzLabelPoints[left](B)
    \tkzLabelPoints[right](C,H)
  \end{tikzpicture}
\end{minipage}
\begin{minipage}{0.33\textwidth}
  \centering
  \begin{tikzpicture}[scale=1]
    \tkzDefPoints{1/2.7/A, 0/0/B, 3.6/0/C}
    \tkzDefPointsBy[projection=onto A--B](C){D}
    \tkzDefPointsBy[projection=onto A--C](B){E}
    \tkzDefPointsBy[projection=onto B--C](A){F}
    \tkzInterLL(A,F)(B,E) \tkzGetPoint{H}

    \tkzDefCircle[circum](C,E,F) \tkzGetPoint{K1}
    \tkzDrawCircles[dashed, thick](K1,C)
    \tkzDefCircle[circum](A,B,E) \tkzGetPoint{K2}
    \tkzDrawCircles[dashed, thick](K2,A)
    \tkzMarkAngles[size=.4, mark=||, mksize=2pt](E,B,A E,F,A E,C,D)
    \tkzFillAngles[size=.4, fill=red!60,opacity=.5](E,B,A E,F,A E,C,D)
    \tkzMarkRightAngles[scale=0.8](A,E,B A,F,B)

    \tkzDrawSegments(A,B A,C B,C B,E A,F C,H)
    \tkzDrawSegments[dashed](H,D E,F)
    \tkzDrawPoints(H)
    % \tkzLabelPoints[above](A)
    \tkzLabelPoints[left](A,B,D)
    \tkzLabelPoints[right](C,H,E)
    \tkzLabelPoints(F)
  \end{tikzpicture}
\end{minipage}
\begin{minipage}{0.33\textwidth}
  \centering
  \begin{tikzpicture}[scale=0.6]
    \tkzDefPoints{1/2.7/A, 0/0/B, 3.6/0/C}
    \tkzDefPointsBy[projection=onto A--B](C){D}
    \tkzDefPointsBy[projection=onto A--C](B){E}
    \tkzDefPointsBy[projection=onto B--C](A){F}
    \tkzInterLL(A,F)(B,E) \tkzGetPoint{H}
    \tkzDefLine[parallel=through A](B,C) \tkzGetPoint{x}
    \tkzDefLine[parallel=through C](A,B) \tkzGetPoint{y}
    \tkzDefLine[parallel=through B](A,C) \tkzGetPoint{z}
    \tkzInterLL(A,x)(C,y) \tkzGetPoint{B'}
    \tkzInterLL(A,x)(B,z) \tkzGetPoint{C'}
    \tkzInterLL(C,y)(B,z) \tkzGetPoint{A'}

    \tkzMarkRightAngles[scale=0.8](C,D,B B,E,C)
    \tkzDrawSegments(A,B A,C B,C B,E C,D A,F)
    \tkzDrawSegments[dashed](A',B' A',C' B',C')
    \tkzDrawPoints(H)
    \tkzLabelPoints[left](B)
    \tkzLabelPoints[right](C,H,A')
    \tkzLabelPoints[below left](A)
    \tkzLabelPoints(C')
    \tkzLabelPoints[below right](B')
  \end{tikzpicture}
\end{minipage}

三条垂线交于一点，证明方法如上图，中图的关键点是$ABFE$四点共圆，用圆周角定理的逆定理判定四点共圆。最右图，作平行线，$H$即为外三角形$A'B'C'$外心，外心必交于一点。

三角形的重心$G$、外心$O$、垂心$H$共线(欧拉线)。

内切圆的半径和三条高的关系$\displaystyle \frac{1}{r}=\frac{1}{h_1}+\frac{1}{h_2}+\frac{1}{h_3}$。面积相等证明。

\subsection{Summary}
一些重要的公式，往往是要解决具体的实际问题，比如：

知道两边和两边的夹角，如何求出第三边，那就是余弦定理

知道三边求角度，显然也是余弦定理。

知道两边和两边的夹角，如何求出三角形面积，那就是$S=\frac{1}{2}ab\sin \theta$。

知道三边长，如何求出三角形面积，这就有了海伦公式。

三角形如何求出外接圆的直径，正弦定理显然合适。

知道三边长，底边一点分底边比例已知，求定点到底边点的距离，显然就是斯图尔特定理。

\clearpage
\section{圆}
《几何原本》中：\\
命题\uppercase\expandafter{\romannumeral3}.20：同弧所对的圆周角是圆心角的一半\\
命题\uppercase\expandafter{\romannumeral3}.21：在同一个圆内，同弧所对的各个圆周角都相等\\
命题\uppercase\expandafter{\romannumeral3}.27：在等圆中，等弧(相等的圆周)上的圆心角或圆周角是相等的\\
命题\uppercase\expandafter{\romannumeral3}.28：在等圆中，相等的直线（弦）截出的弧相等\\
命题\uppercase\expandafter{\romannumeral3}.29：在等圆中，相等的弧对等直线（弦）\\
推论：同圆中或等圆中，同弧对等角，等弧对等角，等角对等弧\\
命题\uppercase\expandafter{\romannumeral3}.31：半圆上的圆周角是直角；大于半圆的圆周角是锐角，小于半圆的圆周角是钝角\\
命题\uppercase\expandafter{\romannumeral3}.35：相交弦定理\\
命题\uppercase\expandafter{\romannumeral3}.36：切割线定理\\
命题\uppercase\expandafter{\romannumeral3}.37：切割线逆定理

命题\uppercase\expandafter{\romannumeral4}.22：圆内接四边形的对角之和等于两直角

事实上，几何原本的命题都建议去自己去证明一下，会很有裨益。

\textbf{四点共圆(Cyclic quadrilateral)判定：}对角互补、圆周角逆定理、托勒密定理、定点到同一平面四点距离相等。核心是对角互补，圆周角的逆定理也可以用相似得出对角互补，外角等于内对角本质也是对角互补，还有相交弦定理，也是可以用相似证明对角互补，托勒密逆定理最后也是对角互补。所以 \textbf{四点共圆的核心是对角互补。}对角互补共圆的判定证明方法，一般用反证法，先三点做出一个圆，假设另外一个点不在圆上，得出矛盾。详见练习。

几何原本中没有现代的“四点共圆”的描述，相关的有命题\uppercase\expandafter{\romannumeral4}.22：圆内接四边形的对角之和等于两直角。命题\uppercase\expandafter{\romannumeral4}.5：可以为任何给定的三角形作外接圆。命题\uppercase\expandafter{\romannumeral3}.25：给定弧可以做出一个圆。

\begin{exercise}
  求证：若凸四边形满足一组对角互补，那么四边形的四个点共圆。
\end{exercise}
\begin{solution}
\begin{minipage}{0.33\textwidth}
  \centering
  \begin{tikzpicture}[scale=1]
    \tkzDefPoints{1/3/A, 0/0/B, 3/0/C}
    \tkzDefCircle[circum](A,B,C) \tkzGetPoints{K}{k}
    \tkzDrawCircles[dashed](K,k)
    \tkzDefPointBy[rotation=center K angle 90](C) \tkzGetPoint{D}
    \tkzDrawSegments(A,B B,C C,D A,D A,C)
    \tkzLabelPoints[above](A)
    \tkzLabelPoints[left](B)
    \tkzLabelPoints[right](C,D)
  \end{tikzpicture}
\end{minipage}
\begin{minipage}{0.33\textwidth}
  \centering
  \begin{tikzpicture}[scale=1]
    \tkzDefPoints{1/3/A, 0/0/B, 3/0/C}
    \tkzDefCircle[circum](A,B,C) \tkzGetPoints{K}{k}
    \tkzDrawCircles[dashed](K,k)
    \tkzDefPointBy[rotation=center K angle 90](C) \tkzGetPoint{D'}
    \tkzDefPointOnLine[pos=0.9](B,D')  \tkzGetPoint{D}
    \tkzInterLC(C,D)(K,k) \tkzGetPoints{C}{E}
    \tkzDrawSegments(A,B B,C C,D A,D A,C)
    \tkzDrawSegments[dashed,red](A,E D,E)
    \tkzLabelPoints[above](A,E)
    \tkzLabelPoints[left](B)
    \tkzLabelPoints[right](C,D)
  \end{tikzpicture}
\end{minipage}
\begin{minipage}{0.33\textwidth}
  \centering
  \begin{tikzpicture}[scale=1]
    \tkzDefPoints{1/3/A, 0/0/B, 3/0/C}
    \tkzDefCircle[circum](A,B,C) \tkzGetPoints{K}{k}
    \tkzDrawCircles[dashed](K,k)
    \tkzDefPointBy[rotation=center K angle 90](C) \tkzGetPoint{D'}
    \tkzDefPointOnLine[pos=1.1](B,D')  \tkzGetPoint{D}
    \tkzInterLC(C,D)(K,k) \tkzGetPoints{C}{E}
    \tkzDrawSegments(A,B B,C C,D A,D A,C)
    \tkzDrawSegments[dashed,red](A,E)
    \tkzLabelPoints[above](A)
    \tkzLabelPoints[left](B)
    \tkzLabelPoints[right](C,D,E)
  \end{tikzpicture}
\end{minipage}

\textbf{【解析】}此题用反证法简单又严谨。上图-左，四边形$ABCD$中$\angle B+\angle D=180^{\circ}$，三角形$ABC$有且唯一的一个外接圆；假设$D$在圆内，如上图-中，$\angle D=\angle DAE+\angle E$，所以$\angle D > \angle E$，$E$在圆上，有$\angle B+\angle E=180^{\circ}$，$\angle B+\angle D>\angle B+\angle E=180^{\circ}$，与已知矛盾。同理可证如果$D$在圆外也不成立，不在圆内也不在圆外，那么只能在圆上，此圆是$\triangle ABC$的外接圆。
\end{solution}

\begin{exercise}
  求证(圆周角逆定理)：若线段$BD$在同侧对点$A,D$的两个角$\angle BAC=\angle BDC$，则$A,B,C,D$四点共圆。\\
\end{exercise}
\begin{solution}
  \begin{minipage}{0.5\textwidth}
    \centering
    \begin{tikzpicture}[scale=1]
      \tkzDefPoints{1/3/A, 0/0/B, 3/0/C}
      \tkzDefCircle[circum](A,B,C) \tkzGetPoints{K}{k}
      \tkzDefPointBy[rotation=center K angle 90](C) \tkzGetPoint{D}
      \tkzDrawSegments(A,B B,C C,D A,C B,D)
      \tkzLabelPoints[above](A)
      \tkzLabelPoints[left](B)
      \tkzLabelPoints[right](C,D)
    \end{tikzpicture}
  \end{minipage}
    \begin{minipage}{0.5\textwidth}
    \centering
    \begin{tikzpicture}[scale=1]
      \tkzDefPoints{1/3/A, 0/0/B, 3/0/C}
      \tkzDefCircle[circum](A,B,C) \tkzGetPoints{K}{k}
      \tkzDefPointBy[rotation=center K angle 90](C) \tkzGetPoint{D}
      \tkzInterLL(A,C)(B,D) \tkzGetPoint{E}
      \tkzLabelAngle[pos=0.5, font=\scriptsize, red](B,A,C){$1$}
      \tkzLabelAngle[pos=0.5, font=\scriptsize, red](B,D,C){$2$}
      \tkzLabelAngle[pos=0.5, font=\scriptsize, red](D,B,A){$3$}
      \tkzLabelAngle[pos=0.5, font=\scriptsize, red](D,C,A){$4$}
      \tkzLabelAngle[pos=0.5, font=\scriptsize, red](C,A,D){$5$}
      \tkzLabelAngle[pos=0.5, font=\scriptsize, red](C,B,D){$6$}
      \tkzLabelAngle[pos=0.5, font=\scriptsize, red](A,D,B){$7$}
      \tkzLabelAngle[pos=0.5, font=\scriptsize, red](A,C,B){$8$}
      \tkzDrawSegments(A,B B,C C,D A,C B,D A,D)
      \tkzLabelPoints[above](A)
      \tkzLabelPoints[left](B,E)
      \tkzLabelPoints[right](C,D)
    \end{tikzpicture}
  \end{minipage}

  \textbf{【解析】}连接$AD$，如上图-右，已知$\angle 1=\angle 2$，那么$\triangle EBA\sim \triangle ECD$，$\angle 3= \angle 4$，$BE:AE=EC:ED$，所以$\triangle EAD \sim \triangle EBC$，进一步得$\angle 5=\angle 6, \; \angle 7 = \angle 8$。$\angle B+\angle D=\angle 3+\angle 6+\angle 7+\angle 2=\angle 3+\angle 5 +\angle 7 +\angle 1=180^{\circ}$。这就是说，四边形$ABCD$，对角互补，那么四点共圆，得证。\\
  Look，圆周角的逆定理的也可指向四边形对角互补。
\end{solution}

\clearpage
\section{直线}
\textbf{一般式：}$Ax+By+C=0$，一般限定$A\ge 0, gcd(A,B,C)=1$

\textbf{两点式：}$\displaystyle \frac{y-y_1}{y_2-y_1}=\frac{x-x_1}{x_2-x_1} \left( \Leftarrow \frac{y-y_1}{x-x_1}=\frac{y_2-y_1}{x_2-x_1}\right)$，此式要求直线不平行于两个坐标轴。\\
\phantom{四个汉字}此式反映了直线上任意点到两定点的斜率相等。

\textbf{斜截式：}$y=mx+b$，斜率为$m$，y轴截距为$b$

\textbf{截距式：}$\displaystyle \frac{x}{a}+\frac{y}{b}=1$，直线与坐标轴的截距分别为$a,b$

\subsection{斜率(slope)}
斜率用来描述直线的方向和陡度，也用来计算斜坡的倾斜程度，通常用$k$表示(不同国家不同的表示，美国、英国等用$m$)。
\[
  k=\frac{\Delta y}{\Delta x} = \frac{\text{rise}}{\text{run}}
\]
对于直角坐标系：$y=kx+b$，若$k$不等于$0$，则可说$k$是斜率

另一个相关概念时倾角(angle of inclination)或斜角，直线与水平轴所夹的最小角。

斜率连接了代数和几何，用来分析“变化快慢”和“倾斜方向”(倾斜程度)。也是直角坐标系中直线的倾斜度的直接表达。

\subsection{点到直线的距离}
若点$P(x_0,y_0)$, 则距离：$\displaystyle  d=\frac{|Ax_0+By_0+C|}{\sqrt{A^2+B^2}}$

证明方法很多，这里用平移坐标系的方式来证明。

我们将新坐标系的原点平移至点$P$，那么$x=x'+x_0, \; y=y'+y_0$，代入，新坐标系内图像的方程就是$A(x'+x_0)+B(y'+y_0)+C=0 \Rightarrow Ax'+By'+Ax_0+By_0+C=0$，与两个轴的交点和原点正好构成一个直角三角形，原点到斜边的距离就是所求，两个直角边分别为$|Ax_0+By_0+C|/|A|, \; |Ax_0+By_0+C|/|B|$, 根据面积相等，距离就是斜边上的高$=$两个直角边的乘积除以斜边长，令$t=|Ax_0+By_0+C|$，简化并代入求解既是。

如何记忆，如果点$P$在直线上，那么距离一定是$0$(分子)，分母可以用直角三角形来记忆，用这两个点理解记忆还是不错的(虽然第二个点有点牵强)。其实直线的法向量是$(A,B)$，分母就是法向量的长度。关于记忆，用一个简单的例子来辅助也是不错的，比如原点到直线$x+y=1$的距离。

\subsection{平行线之间的距离}
假设两条直线分别为$Ax+By+C_1=0, \; Ax+By+C_2=0$，假设$P(x_0,y_0)$在第一条在直线上，那么$P$到另外一条直线上的距离就是:$\displaystyle  d=\frac{|Ax_0+By_0+C_2|}{\sqrt{A^2+B^2}}$,
while $Ax_0+By_0+C_1=0 \Rightarrow Ax_0+By_0=-C_1$代入，得：$\displaystyle d=\frac{|C_1-C_2|}{\sqrt{A^2+B^2}}$

\clearpage
\section{一元二次方程}

$ax^2+bx+c=0$, 两个根$x_1, x_2$, $\Delta = b^2-4ac$, 对称轴$-\frac{b}{2a}$

$x_{1,2}=\frac{-b\pm \sqrt{b^{2}-4ac}}{2a}, \;\;x_1+x_2=-b/a$, \; $x_1x_2=c/a$

对称轴$-b/2a$, 两个根与对称轴的偏移量就是$|\sqrt{b^2-4ac}/2a|$，注意到两根是对称的，有时候可以把两个根设为$m-n, m+n$，其中$m$是对称轴。比如$x^2-8x+12=0$，Let $x_1=4-n, x_2=4+n$，根据韦达，
$x_1x_2=12=(4-n)(4+n)$，很容易算出偏移量。

通过变型可得$ax^2=-bx-c$，一个二次型等于一个一次型，emmm...?

\subsection{例子}
例： $x^2-2ax+b=0$有两个根$x_1, x_2$, $y^2+2ax+b=0$有两个根$y_1, y_2$, and $a, b \in \mathbb{Z}^+, x_1y_2-x_2y_1=104$，求$b$的最小值。

已知$a,b$满足$a^2-15a-5=0, b^2-15b-5=0$，求$a/b+b/a$, tips\footnote{$a,b$显然是$x^2-15x-5=0$的两个根}

已知$3(a-b)+\sqrt{3}(b-c)+(c-a)=0$，求$\frac{b-a}{b-c}$, tips\footnote{观察系数$3, \sqrt{3}$，如果$1$代替$\sqrt{3}$显然也成立，所以可知$1, \sqrt{3}$是方程$(a-b)x^2+(b-c)x+(c-a)$的两个根}

已知$x^5+x+1=0$，求$x^3-x^2$, tips\footnote{思路1：已知一个多项式的值，求另一个多项式，可以用长除法，高次的除以低次的，$(x^5+x+1)/(x^3-x^2)=x^2+x+1+R(x^2+x+1) \Rightarrow x^5+x+1=(x^3-x^2)(x^2+x+1)+x^2+x+1$\\
  思路2：Let$x^3-x^2=t$，升次降次到已知的多项式，$x^3=x^2+t, x^4=x^3+tx=x^2+tx+t, x^5=(1+t)x^2+tx+t \Rightarrow (1+t)x^2+(1+t)x+(1+t)=0$
}
\clearpage
\section{几何模型}
\subsection{费马点(Fermat point)}
三角形内(包括边线)一点，到三个顶点的和最小。

\begin{minipage}{0.3\textwidth}
  \centering
  \begin{tikzpicture}
    \tkzDefPoints{0/0/B,3/0/C,2/2.5/A, 1.5/1/P}
    \tkzDrawSegments(A,B B,C A,C P,A P,B P,C)
    \tkzLabelPoints(B,C,P)
    \tkzLabelPoints[above](A)
  \end{tikzpicture}
\end{minipage}
\begin{minipage}{0.3\textwidth}
  \begin{tikzpicture}
    \tkzDefPoints{0/0/B,3/0/C,2/2.5/A, 1.5/1/P}
    \tkzDrawSegments(A,B B,C A,C P,A P,B P,C)

    \tkzDefPointBy[rotation=center B angle 60](A) \tkzGetPoint{A'}
    \tkzDefPointBy[rotation=center B angle 60](P) \tkzGetPoint{P'}
    \tkzDrawSegments[dashed](P',B P',A' A',B A,A' P,P')

    \tkzFillPolygon[color=orange!30, opacity=0.3](A,B,P)
    \tkzFillPolygon[color=orange!30, opacity=0.3](A',B,P')
    \tkzMarkSegments[mark=||, size=2pt](P',B P,P' P,B)

    \tkzMarkAngles[scale=0.6, ->](P,B,P')
    \tkzFillAngles[scale=0.6, fill=red!20,opacity=.5](P,B,P')
    \tkzLabelAngle[pos=0.9, red, scale=0.6](P,B,P'){$60^{\circ}$}

    \tkzLabelPoints(B,C,P)
    \tkzLabelPoints[above](A)
    \tkzLabelPoints[above right](P')
    \tkzLabelPoints[left](A')
  \end{tikzpicture}
\end{minipage}
\begin{minipage}{0.4\textwidth}
\centering
  \begin{tikzpicture}[scale=0.6]
    \tkzDefPoints{0/0/B,3/0/C,2/2.5/A}
    % \tkzDrawPoints(A,B,C)
    \tkzDrawSegments(A,B B,C A,C)

    \tkzDefPointBy[rotation=center B angle 60](A) \tkzGetPoint{C'}
    \tkzDefPointBy[rotation=center B angle -60](C) \tkzGetPoint{A'}
    \tkzDefPointBy[rotation=center C angle -60](A) \tkzGetPoint{B'}
    \tkzDrawSegments[dashed](B,C' A,C' B,A' C,A' C,B' B',A C,C' A,A' B,B')

    \tkzInterLL(A,A')(B,B') \tkzGetPoint{P}

    \tkzLabelPoints(A')
    \tkzLabelPoints[above](A)
    \tkzLabelPoints[left](B, C')
    \tkzLabelPoints[right](C, B', P)
  \end{tikzpicture}
\end{minipage}

三角形三个边分别向外作等边三角形，顶点与对边的连线相交的点就是所有点。其实，三角形要看是否有角度大于$120^{\circ}$，当有一个角不小于$120^{\circ}$时，那么此时费马点就是此顶点。也可以看到，费马点在内部时，费马点和任意点形成的夹角都是$120^{\circ}$。顶角超过$120^{\circ}$，旋转后的$A'C$已越过$A$。

物理学解释：平面上所给的三个给定点钻出洞来，再设有三条绳子系在一起，每条绳子各穿过一个洞口，而绳子的末端都绑有一个固定重量m的重物。假设摩擦力可以忽略，那么绳子会被拉紧，而绳结最后会停在平面一点的上方。

\textbf{加权费马点}，也就是带系数的距离和，比如$PA+\sqrt{2}PB+PC$，或者$PC+\frac{\sqrt{3}}{2}PB+\frac{1}{2}PA$

\begin{minipage}{0.4\textwidth}
  \centering
  \begin{tikzpicture}
    \tkzDefPoints{0/0/B,3/0/C,2/2.5/A, 1.5/1/P}
    \tkzDrawSegments(A,B B,C A,C P,A P,B P,C)

    \tkzDefPointBy[rotation=center B angle 90](A) \tkzGetPoint{A'}
    \tkzDefPointBy[rotation=center B angle 90](P) \tkzGetPoint{P'}
    \tkzDrawSegments[dashed](P',B P',A' A',B A,A' P,P')

    \tkzFillPolygon[color=orange!30, opacity=0.3](A,B,P)
    \tkzFillPolygon[color=orange!30, opacity=0.3](A',B,P')

    \tkzMarkAngles[scale=0.6, ->](P,B,P')
    \tkzFillAngles[scale=0.6, fill=red!20,opacity=.5](P,B,P')
    % \tkzLabelAngle[pos=0.9, red, scale=0.6](P,B,P'){$60^{\circ}$}

    \tkzLabelPoints(B,C,P)
    \tkzLabelPoints[above](A)
    \tkzLabelPoints[above right](P')
    \tkzLabelPoints[left](A')
  \end{tikzpicture}
\end{minipage}
\begin{minipage}{0.3\textwidth}
  \centering
  \begin{tikzpicture}
    \tkzDefPoints{0/0/B,3/0/C,2/2.5/A, 1.5/1/P}
    \tkzDrawSegments(A,B B,C A,C P,A P,B P,C)

    \tkzDefPointBy[rotation=center B angle 60](A) \tkzGetPoint{A'}
    \tkzDefPointBy[rotation=center B angle 60](P) \tkzGetPoint{P'}
    \tkzDrawSegments[dashed](P',B P',A' A',B)

    \tkzFillPolygon[color=orange!30, opacity=0.3](A,B,P)
    \tkzFillPolygon[color=orange!30, opacity=0.3](A',B,P')

    \tkzMarkAngles[scale=0.6, ->](P,B,P')
    \tkzFillAngles[scale=0.6, fill=red!20,opacity=.5](P,B,P')

    \tkzDefPointOnLine[pos=0.5](B,A')  \tkzGetPoint{D}
    \tkzDefPointOnLine[pos=0.5](B,P')  \tkzGetPoint{E}
    \tkzDrawSegments[dashed](D,E E,P)

    \tkzLabelPoints(B,C,P)
    \tkzLabelPoints[above](A)
    \tkzLabelPoints[above right](P', E)
    \tkzLabelPoints[left](A', D)
  \end{tikzpicture}
\end{minipage}
\begin{minipage}{0.3\textwidth}
  \centering
  \begin{tikzpicture}
    \tkzDefPoints{0/0/B,3/0/C,2/2.5/A}
    \tkzDrawSegments(A,B B,C A,C)

    \tkzDefPointBy[rotation=center B angle 60](A) \tkzGetPoint{A'}
    \tkzDrawSegments[dashed](A',B)

    \tkzMarkAngles[scale=0.6, ->](A,B,A')
    \tkzFillAngles[scale=0.6, fill=red!20,opacity=.5](A,B,A')

    \tkzDefPointOnLine[pos=0.5](B,A')  \tkzGetPoint{D}
    \tkzDefPointsBy[projection=onto C--D](B){E}
    \tkzDefPointOnLine[pos=2](B,E)  \tkzGetPoint{P'}

    \tkzDefPointBy[rotation=center B angle -60](P') \tkzGetPoint{P}
    \tkzDrawSegments[dashed](B,P)

    \tkzDrawSegments[dashed](D,C B,P')

    \tkzLabelPoints(B,C)
    \tkzLabelPoints[above](A, P', P, E)
    \tkzLabelPoints[left](A', D)
  \end{tikzpicture}
\end{minipage}

对于有两个不是$1$系数的情况，刚开始的时候，我的思路是：旋转后的$PP'$的范围是$(0,2PB)$，而$PA$可以任意系数，延长或者截取$P'A'$即可。可是本质上的点到点的距离直线段最短，$P$是不定点，延长或截取$P'A'$后的点也是不定点，所以这个思路直接cut掉。这个定点只能在固定的边线$A'B$上取。所以，加权费马点的核心步骤就是：\textbf{1、先找到系数为$1$的边不动, 2、上面的边的系数用平行线，3、另一个系数(旋转点的边)用旋转，用勾股定理确定角度}。如上图，系数$PA$是用平行线转化为$DE$，$PE$是带系数的$PB$转化的(因为$BE$是截取的$PB$，$\triangle PBE$的三边就确定了，旋转的角度$\angle PBE$可用勾股定理计算出)，就转化为$C,D$两点之间线段最短($CD=DE+PE+PC$)。其实，有时候很反直观，比如上中图，平行线$DE$的向下很陡，如何到$C$呢？那又如何确定这个点呢？

要确定点就需要反推。例如上又图，$D$点是固定的，旋转的角度是固定的，连接$CD$，那这个$P$点肯定在之上，$\angle BEC$是定的($\triangle PBE$三边已知，一般很有可能是特殊角，比如垂直)，所以能找到$E$点，延长找到$P'$点，反旋转后找到P点。还有一种方法就是，往右旋转，找到$B', D'$点，$CD$和$BD'$的连线交点即是。

思考：$10PA+2PB+2PC$，如果往外作平行线会不会定点过高，然后两个定点的连线会越过三角形？是不是最好往小了缩放？

\subsection{探照灯模型- 定角定高}

\begin{minipage}[c]{0.33\textwidth}
  \centering
  \begin{tikzpicture}
    \tkzDefPoints{0/0/A,3/0/B, 2.1/3/C}
    \tkzDrawPoints(A,B,C)
    \tkzDrawSegments(A,C B,C)
    \tkzDrawLines[red, thick](A,B)

    \tkzMarkAngles[scale=0.8](A,C,B)
    \tkzFillAngles[scale=0.8, fill=red!20,opacity=.5](A,C,B)
    \tkzLabelAngle[pos=0.9, red](A,C,B){$\theta$}

    \tkzDefPointsBy[projection=onto A--B](C){D}
    \tkzDrawSegment[thick, red](C,D)
    \tkzLabelSegment[right, midway, red](C,D){h}
    \tkzMarkRightAngle(A,D,C)
    \tkzLabelPoints(A,B,D)
    \tkzLabelPoints[right](C)
  \end{tikzpicture}
\end{minipage}
\begin{minipage}[c]{0.33\textwidth}
  \centering
  \begin{tikzpicture}
    \tkzDefPoints{0/0/A,3/0/B, 2.1/3/C}
    \tkzDrawPoints(A,B,C)
    \tkzDrawSegments(A,C B,C)
    \tkzDrawLines[red, thick](A,B)

    \tkzDefCircle[circum](A,B,C) \tkzGetPoint{O}
    \tkzDrawCircle[dashed, thick](O,A)

    \tkzDrawSegments[dashed](O,A O,B O,C)

    \tkzMarkAngles[scale=0.8](A,C,B)
    \tkzFillAngles[scale=0.8, fill=red!20,opacity=.5](A,C,B)
    \tkzLabelAngle[pos=0.9, red](A,C,B){$\theta$}

    \tkzDefPointsBy[projection=onto A--B](C){D}
    \tkzDrawSegment[thick, red](C,D)
    \tkzLabelSegment[right, midway, red](C,D){h}
    \tkzMarkRightAngle(A,D,C)

    \tkzDefPointsBy[projection=onto A--B](O){E}
    \tkzDrawSegment[dashed](O,E)
    \tkzMarkRightAngle[scale=0.8](A,E,O)

    \tkzMarkAngles[scale=0.4](A,O,E)
    \tkzMarkAngles[scale=0.45](E,O,B)
    \tkzFillAngles[scale=0.7, fill=red!20,opacity=.5](A,O,E)
    \tkzFillAngles[scale=0.8, fill=red!20,opacity=.5](E,O,B)
    \tkzLabelAngle[pos=0.7](A,O,E){$\theta$}
    \tkzLabelAngle[pos=0.7](E,O,B){$\theta$}

    \tkzLabelPoints(A,B,D,E)
    \tkzLabelPoints[above left](O)
    \tkzLabelPoints[right](C)
  \end{tikzpicture}
\end{minipage}
\begin{minipage}[c]{0.33\textwidth}
  \centering
  \begin{tikzpicture}
    \tkzDefPoints{0/0/A,3/0/B, 2.1/0.8/C}
    \tkzDrawPoints(A,B,C)
    \tkzDrawSegments(A,C B,C)
    \tkzDrawLines[red, thick](A,B)

    \tkzDefCircle[circum](A,B,C) \tkzGetPoint{O}
    \tkzDrawCircle[dashed, thick](O,A)
    \tkzLabelPoints(O)

    \tkzDrawSegments[dashed](O,A O,B O,C)

    \tkzMarkAngles[scale=0.6](A,C,B)
    \tkzFillAngles[scale=0.6, fill=red!20,opacity=.5](A,C,B)

    \tkzDefPointsBy[projection=onto A--B](C){D}
    \tkzDrawSegment[thick, red](C,D)
    \tkzLabelSegment[right, midway, red](C,D){h}
    \tkzMarkRightAngle[scale=0.8](A,D,C)

    \tkzDefPointsBy[projection=onto A--B](O){E}
    \tkzDrawSegment[dashed](O,E)
    \tkzMarkRightAngle[scale=0.8](A,E,O)

    \tkzMarkAngles[scale=0.6](E,O,A)
    \tkzFillAngles[scale=0.6, fill=red!20,opacity=.5](E,O,A)
    \tkzLabelAngle[pos=0.4, red, scale=0.8](E,O,A){$90^{\circ}-\frac{\theta}{2}$}
    \tkzDrawPoints(O)

    \tkzLabelPoints(A,B,D)
    \tkzLabelPoints[right](C)
    \tkzLabelPoints[above](E)
  \end{tikzpicture}
\end{minipage}

定点到定直线的距离是定值，顶角是定值, 求三角形底边最小值(三角形面积最小)。

$AB=2r\sin\theta$，只需找出最小外接圆半径即可， while $h\le r+r\cos\theta \Rightarrow r \ge h/(1+\cos\theta)$ or $h+r\cos(90^{\circ}-\theta/2) \le r \Rightarrow r\ge h/(1-\sin \frac{\theta}{2})$。最小值的情况下，都是$0C \perp AB$，也就是$OC$平分定角，也就是$CA=CB$，$\triangle CAB$是等腰三角形。

现在我们换一种思路，用代数的方法，Let $\angle ACD=\angle\alpha$，有$AB=h\tan\alpha+h\tan(\theta-\alpha)$，可以转化为$f(\alpha)=h\left[\tan\alpha+\tan(\theta-\alpha)\right]$的最值问题，求导$f'(\alpha)=h \frac{1}{\cos^{2}\alpha}-h \frac{1}{\cos^2(\theta-\alpha)}=0 \;\Rightarrow\cos^2\alpha=\cos^2(\theta-\alpha)$。默认的条件是$0<\alpha<\theta<\pi$，分类讨论，$\cos\alpha=\cos(\theta-\alpha)\Rightarrow \cos\alpha - \cos(\theta-\alpha)=0$，利用和差化积得$2\sin(\theta/2)\sin(\frac{2\alpha-\theta}{2})=0 \Rightarrow 2\alpha=\theta$；或者$\cos\alpha=-\cos(\theta-\alpha)$，不符合。\todo{???}利用求导求最值还需要比较端点和驻点。

\clearpage
\section{Exercises}

\begin{exercise}
  \begin{minipage}{0.3\textwidth}
    \centering
    \begin{tikzpicture}[scale=1]
      \tkzDefPoints{0/0.2/A, .1/1.5/B, 3/1.9/C}
      \tkzDefCircle[circum](A,B,C) \tkzGetPoints{K}{k}
      \tkzDrawCircles(K,k)
      \tkzDefMidArc(K,C,A) \tkzGetPoint{M}
      \tkzDefMidArc(K,A,C) \tkzGetPoint{M'}
      \tkzDefPointsBy[projection=onto B--C](M,M'){H,H'}
      \tkzMarkRightAngles[scale=0.7](M,H,C M',H',C)
      \tkzDrawSegments(A,B B,C M,H M',H')
      \tkzDrawPoints(K)
      \tkzLabelPoints(H)
      \tkzLabelPoints[above](H')
      \tkzLabelPoints[left](A,B)
      \tkzLabelPoints[right](C,M')
      \tkzLabelPoints[above left](M)
    \end{tikzpicture}
  \end{minipage}
  \begin{minipage}{0.7\textwidth}
    如图：$AB,BC$是圆的两条弦，$BC>AB$，$M,M'$是弧$AC$的两个中点，$MH\perp BC, \;M'H'\perp BC$，\\
    求证：$CH=AB+BH, \; CH'=BH'-AB$
  \end{minipage}
  \end{exercise}
\begin{solution}
  \textbf{【解析】}这就是阿基米德折弦(中点)定理(Archimedes' Midpoint Theorem)，证明方法有补短法、截长法、垂线法等。用全等，都需要内接圆对角和性质和等弧对等角。

  \begin{minipage}{0.4\textwidth}
    \centering
    \begin{tikzpicture}[scale=1]
      \tkzDefPoints{0/0.2/A, .1/1.5/B, 3/1.9/C}
      \tkzDefCircle[circum](A,B,C) \tkzGetPoints{K}{k}
      \tkzDrawCircles(K,k)
      \tkzDefMidArc(K,C,A) \tkzGetPoint{M}
      \tkzDefMidArc(K,A,C) \tkzGetPoint{M'}
      \tkzDefPointsBy[projection=onto B--C](M,M'){H,H'}
      \tkzMarkRightAngles[scale=0.7](M,H,C)
      \tkzDefPointsBy[symmetry=center H](C){D}

      \tkzDrawSegments(A,B B,C M,H)
      \tkzDrawSegments[dashed](M,D B,D M,C M,B A,C M,A)
      \tkzDrawPoints(K)
      \tkzLabelPoints(H,D)
      \tkzLabelPoints[left](A)
      \tkzLabelPoints[below left](B)
      \tkzLabelPoints[right](C)
      \tkzLabelPoints[above left](M)
    \end{tikzpicture}
  \end{minipage}
  \begin{minipage}{0.3\textwidth}
    \centering
    \begin{tikzpicture}[scale=1]
      \tkzDefPoints{0/0.2/A, .1/1.5/B, 3/1.9/C}
      \tkzDefCircle[circum](A,B,C) \tkzGetPoints{K}{k}
      \tkzDrawCircles(K,k)
      \tkzDefMidArc(K,C,A) \tkzGetPoint{M}
      \tkzDefPointsBy[projection=onto B--C](M){H}
      \tkzMarkRightAngles[scale=0.7](M,H,C)
      \tkzDefPointsBy[symmetry=center H](B){B'}

      \tkzDrawSegments(A,B B,C M,H)
      \tkzDrawSegments[dashed](M,A M,C A,C M,B M,B')
      \tkzDrawPoints(K)
      \tkzLabelPoints(H,B')
      \tkzLabelPoints[left](A,B)
      \tkzLabelPoints[right](C)
      \tkzLabelPoints[above left](M)
    \end{tikzpicture}
  \end{minipage}
  \begin{minipage}{0.3\textwidth}
    \centering
    \begin{tikzpicture}[scale=1]
      \tkzDefPoints{0/0.2/A, .1/1.5/B, 3/1.9/C}
      \tkzDefCircle[circum](A,B,C) \tkzGetPoints{K}{k}
      \tkzDrawCircles(K,k)
      \tkzDefMidArc(K,C,A) \tkzGetPoint{M}
      \tkzDefMidArc(K,A,C) \tkzGetPoint{M'}
      \tkzDefPointsBy[projection=onto B--C](M,M'){H,H'}
      \tkzDefPointsBy[projection=onto A--B](M){D}
      \tkzMarkRightAngles[scale=0.7](M,H,C B,D,M)

      \tkzDrawSegments(A,B B,C M,H)
      \tkzDrawSegments[dashed](M,D B,D M,A M,C M,B A,C)
      \tkzDrawPoints(K)
      \tkzLabelPoints(H)
      \tkzLabelPoints[left](A,B,D)
      \tkzLabelPoints[right](C)
      \tkzLabelPoints[above left](M)
    \end{tikzpicture}
  \end{minipage}
\end{solution}

\hrulefill

\begin{exercise}
  \begin{minipage}{0.3\textwidth}
    \centering
    \begin{tikzpicture}[scale=1]
      \tkzDefPoints{.7/3.3/A, 0/0/B, 2.4/0/C, 1/3.3/A'}
      \tkzDefLine[bisector](A,B,C)  \tkzGetPoint{b}
      \tkzDefPointOnLine[pos=0.5](A,B)  \tkzGetPoint{F}
      \tkzDefLine[parallel=through F](B,C) \tkzGetPoint{F'}
      \tkzInterLL(F,F')(B,b) \tkzGetPoint{E}
      \tkzInterLL(C,E)(A,A') \tkzGetPoint{D}
      \tkzDrawPolygons(A,B,C,D)
      \tkzDrawSegments(A,B A,E B,E)
      \tkzLabelPoints[left](A,B)
      \tkzLabelPoints[right](C,E,D)
    \end{tikzpicture}
  \end{minipage}
  \begin{minipage}{0.7\textwidth}
    梯形$ABCD$，$AD \parallel BC$，$AB=AD+BC$。\\
    求证：$\angle A, \angle B$的角平分线比相交于另一腰上。
  \end{minipage}
\end{exercise}
\begin{solution}
  此题需要先进行探讨，权且命题成立，去打听这个腰上的点具体的“消息”。如下左图，很明显全等得出$DE=EF=EC$，也就是$E$是腰的中点。那我们就顺着这个思路，反过去证。

  先做$\angle B$的角平分线，交$CD$于$E$点，连接$AE$，如果$AE$平分$\angle A$，那么就得证。$\triangle BFE \cong \triangle BCE$很容易，接下来要求证$\angle 1=\angle 2$。$\triangle AEF \cong \triangle AED$，限于现在的辅助线，$AF=AD, AE=AE$，条件不足了。这时候就分析一下，\textbf{什么条件没用上？} 很明显，是已知的梯形，梯形的上下边平行。这时候就应该需要额外的辅助线了。我们就延长$BE$，如下中图，$\angle 4 = \angle G$，$\triangle ABE$是等腰，$\angle 3 = \angle G$，可以推出$AB=AG$，瞧，我们可以得出$DG=BC$，很明显$\triangle EDG \cong \triangle ECB$。局势瞬间就明朗了，之前的辅助点$F$也就不需要了。

  可是，总觉得，哪里有点“不舒服”，我们设想，如果$E$为腰的中点，那么$AE,BE$是不是必平分两个角呢？很显然也是。而且梯形的中位线长度正好等于$\frac{1}{2}(BC+AD)$，可以在梯形的中位线做点文章，如下图右，看似更简单了不是？

  \textbf{中位线定理在三角形和梯形题型中的重要性也是独一无二的存在。}

  \begin{minipage}{0.3\textwidth}
    \centering
    \begin{tikzpicture}[scale=1]
      \tkzDefPoints{.7/3.3/A, 0/0/B, 2.4/0/C, 1/3.3/A'}
      \tkzDefLine[bisector](A,B,C)  \tkzGetPoint{b}
      \tkzDefPointOnLine[pos=0.5](A,B)  \tkzGetPoint{F}
      \tkzDefLine[parallel=through F](B,C) \tkzGetPoint{F'}
      \tkzInterLL(F,F')(B,b) \tkzGetPoint{E}
      \tkzInterLL(C,E)(A,A') \tkzGetPoint{D}
      \tkzDefPointsBy[reflection=over B--E](C){F}

      \tkzLabelAngle[pos=0.5, font=\scriptsize, red](E,A,D){$1$}
      \tkzLabelAngle[pos=0.5, font=\scriptsize, red](F,A,E){$2$}
      \tkzLabelAngle[pos=0.5, font=\scriptsize, red](E,B,A){$3$}
      \tkzLabelAngle[pos=0.5, font=\scriptsize, red](C,B,E){$4$}
      \tkzDrawPolygons(A,B,C,D)
      \tkzDrawSegments(A,B A,E B,E)
      \tkzDrawSegments[dashed](E,F)
      \tkzLabelPoints[left](A,B,F)
      \tkzLabelPoints[right](C,E,D)
    \end{tikzpicture}
  \end{minipage}
  \begin{minipage}{0.4\textwidth}
    \centering
    \begin{tikzpicture}
      \tkzDefPoints{.7/3.3/A, 0/0/B, 2.4/0/C, 1/3.3/A'}
      \tkzDefLine[bisector](A,B,C)  \tkzGetPoint{b}
      \tkzDefPointOnLine[pos=0.5](A,B)  \tkzGetPoint{F}
      \tkzDefLine[parallel=through F](B,C) \tkzGetPoint{F'}
      \tkzInterLL(F,F')(B,b) \tkzGetPoint{E}
      \tkzInterLL(C,E)(A,A') \tkzGetPoint{D}
      \tkzInterLL(B,E)(A,D) \tkzGetPoint{G}

      \tkzLabelAngle[pos=0.5, font=\scriptsize, red](E,A,D){$1$}
      \tkzLabelAngle[pos=0.5, font=\scriptsize, red](F,A,E){$2$}
      \tkzLabelAngle[pos=0.5, font=\scriptsize, red](E,B,A){$3$}
      \tkzLabelAngle[pos=0.5, font=\scriptsize, red](C,B,E){$4$}
      \tkzDrawPolygons(A,B,C,D)
      \tkzDrawSegments(A,B A,E B,E)
      \tkzDrawSegments[dashed](E,G A,G)
      \tkzLabelPoints[left](A,B)
      \tkzLabelPoints[right](C,E,D,G)
    \end{tikzpicture}
  \end{minipage}
  \begin{minipage}{0.3\textwidth}
    \centering
    \begin{tikzpicture}
      \tkzDefPoints{.7/3.3/A, 0/0/B, 2.4/0/C, 1/3.3/A'}
      \tkzDefLine[bisector](A,B,C)  \tkzGetPoint{b}
      \tkzDefPointOnLine[pos=0.5](A,B)  \tkzGetPoint{F}
      \tkzDefLine[parallel=through F](B,C) \tkzGetPoint{F'}
      \tkzInterLL(F,F')(B,b) \tkzGetPoint{E}
      \tkzInterLL(C,E)(A,A') \tkzGetPoint{D}
      \tkzDefPointOnLine[pos=0.5](A,B)  \tkzGetPoint{F}

      \tkzLabelAngle[pos=0.5, font=\scriptsize, red](E,A,D){$1$}
      \tkzLabelAngle[pos=0.5, font=\scriptsize, red](F,A,E){$2$}
      \tkzLabelAngle[pos=0.5, font=\scriptsize, red](E,B,A){$3$}
      \tkzLabelAngle[pos=0.5, font=\scriptsize, red](C,B,E){$4$}
      \tkzLabelAngle[pos=0.5, font=\scriptsize, red](A,E,F){$5$}
      \tkzLabelAngle[pos=0.5, font=\scriptsize, red](F,E,B){$6$}
      \tkzDrawPolygons(A,B,C,D)
      \tkzDrawSegments(A,B A,E B,E)
      \tkzDrawSegments[dashed](E,F)
      \tkzLabelPoints[left](A,B,F)
      \tkzLabelPoints[right](C,E,D)
    \end{tikzpicture}
  \end{minipage}
\end{solution}

\hrulefill

\begin{exercise}
  \begin{minipage}{0.3\textwidth}
    \centering
    \begin{tikzpicture}
      \tkzDefPoints{1.6/2.7/A, 0/0/B, 3/0.2/C}

      \tkzDefCircle[circum](A,B,C) \tkzGetPoint{O}
      \tkzDrawCircles[thick](O,A)
      \tkzDefPointsBy[projection=onto B--C](A){D}
      \tkzDefPointsBy[projection=onto A--B](D){E}
      \tkzDefPointsBy[projection=onto A--C](D){F}
      \tkzMarkRightAngles[scale=0.8](D,E,B D,F,C A,D,B)

      \tkzDrawSegments(A,B A,C B,C A,D D,E D,F E,F)
      \tkzDrawPoints(O)
      \tkzLabelPoints[above](A,O)
      \tkzLabelPoints(D)
      \tkzLabelPoints[left](B,E)
      \tkzLabelPoints[right](C,F)
    \end{tikzpicture}
  \end{minipage}
  \begin{minipage}{0.7\textwidth}
    $\triangle ABC$的外接圆半径为$R$, $AD\perp BC, \; DE\perp AB, \; DF\perp AC$。\\求证：$S_{\triangle ABC}=R\cdot EF$.\\
  \end{minipage}
\end{exercise}
\begin{solution}
  \begin{minipage}{0.7\textwidth}
    \textbf{【解析】}如果觉得用外接圆证明($AEDF$共圆，$AD$是直径)有点索然无趣，那就用作辅助线找图形面积相等吧，还会得到一个意外的收获($AH \perp EF$)
  \end{minipage}
  \begin{minipage}{0.3\textwidth}
    \centering
    \begin{tikzpicture}
      \tkzDefPoints{1.8/2.7/A, 0/0/B, 3/0.2/C}

      \tkzDefCircle[circum](A,B,C) \tkzGetPoint{O}
      \tkzDrawCircles[thick](O,A)
      \tkzDefPointsBy[projection=onto B--C](A){D}
      \tkzDefPointsBy[projection=onto A--B](D){E}
      \tkzDefPointsBy[projection=onto A--C](D){F}
      \tkzMarkRightAngles[scale=0.8](D,E,B D,F,C A,D,B)
      \tkzInterLC(A,O)(O,A) \tkzGetPoints{H}{A}

      \tkzDrawSegments(A,B A,C B,C A,D D,E D,F E,F)
      \tkzDrawSegments[dashed](A,H H,E H,F H,B H,C)
      \tkzDrawPoints(O)
      \tkzLabelPoints[above](A)
      \tkzLabelPoints(D)
      \tkzLabelPoints[left](B,E,H)
      \tkzLabelPoints[right](C,O)
      \tkzLabelPoints[above right](F)
    \end{tikzpicture}
  \end{minipage}
\end{solution}

\hrulefill

\begin{exercise}
  \begin{minipage}{0.3\textwidth}
    \centering
    \begin{tikzpicture}[scale=0.8]
      \tkzDefPoints{0/4/A, 0/0/B, 3/0/C, 0.6/4/D}
      \tkzDrawPolygons(A,B,C,D)
      \tkzDefPointOnLine[pos=0.5](C,D)  \tkzGetPoint{O}
      \tkzDrawSemiCircle(O,D)
      \tkzInterLC(A,B)(O,C) \tkzGetPoints{E}{F}
      \tkzMarkRightAngles[scale=0.8](B,A,D A,B,C)

      \tkzDrawPoints(O)
      \tkzLabelPoints[left](A,B)
      \tkzLabelPoints[right](C,D,E,F,O)
    \end{tikzpicture}
  \end{minipage}
  \begin{minipage}{0.7\textwidth}
    如左图，直角梯形中，$\angle A=\angle B=90^{\circ}$，$AB=a, \; BC=b, \; AD=c$，以$CD$为直径的圆与$AB$的交于$E,F$。\\
    求证：$AE,BE$是方程$x^2-ax+bc=0$的两个根。
  \end{minipage}
\end{exercise}
\begin{solution}
  \textbf{【解析】}根据韦达定理，很明显$AE+BE=AB=a$，那么要证明$AE\cdot BE=bc$，$AE$和$BE$相乘，而$AB$是直线与圆的相交，应该考的是割线定理，但是位置有点不太对，应该是$AE\cdot AF$才对，所以，需要求证$AF$和$BE$之间的关系，而$AE$和$BF$看起来是相等的，如果$AE=BF$，那么$BE=AF$，就是可以用割线定理。而求证的后面有个$c$，也就是$AD$，也是与圆相交，所以考虑延长一下$AD$，这样辅助线就有了，那就应该是割线定理没跑了。
\end{solution}

\hrulefill

\begin{exercise}
  \begin{minipage}{0.3\textwidth}
    \centering
    \begin{tikzpicture}[scale=1]
      \tkzDefPoints{2.7/2/A, 0/0/B, 3/-0.2/C, 1.5/-0.1/K}
      \tkzDefCircle[circum](A,B,C) \tkzGetPoint{O}
      \tkzDrawCircles(O,A)

      \tkzInterLC(O,K)(O,A) \tkzGetPoints{M'}{M}
      \tkzDefPointsBy[projection=onto A--B](M'){D}
      \tkzDefPointsBy[projection=onto C--A](M'){E}
      \tkzMarkRightAngles[scale=0.8](M',D,A M',E,A)

      \tkzDrawSegments(A,B A,C B,C  M',M M',D M',E A,E)
      \tkzDrawPoints(O)
      \tkzLabelPoints(M,D)
      \tkzLabelPoints[above](M',E)
      \tkzLabelPoints[left](B)
      \tkzLabelPoints[right](A,C,O)
    \end{tikzpicture}
  \end{minipage}
  \begin{minipage}{0.7\textwidth}
    $\triangle ABC$内接于圆$O$，$M$为$\overset{\frown}{BC}$的中点，$MM'$为直径，\\
    $M'D\perp AB, \; M'E\perp AC$。\\求证：$DA+EA=AB-AC$
  \end{minipage}
\end{exercise}
\begin{solution}
  \textbf{【解析】}此题明显也是追溯型的题型，几何中$DA+EA$和$AB-AC$无从下手，移项$AC+EA=AB-DA$，正好是$CE=BD$，又看到垂直，连接$BM',M'C$，直角三角形，已知可知$M'B=M'C$，所以只要再找到一个角即可，很明显$\angle M'BD = \angle M'CE$(圆周角定理)。
\end{solution}

\hrulefill

\begin{exercise}
  \begin{minipage}{0.3\textwidth}
    \centering
    \begin{tikzpicture}[scale=1.2]
      \tkzDefPoints{0/1/A, 0/0/B, 3/0/C}
      \tkzDefPointBy[rotation=center A angle 45](C) \tkzGetPoint{C'}
      \tkzDefPointOnLine[pos=0.5](A,C')  \tkzGetPoint{D}

      \tkzDrawSegments(A,B A,C B,C A,D B,D C,D)
      \tkzLabelPoints[above](A)
      \tkzLabelPoints[right](D)
      \tkzLabelPoints[](B,C)

      \tkzMarkAngles[scale=0.5](C,A,D)
      \tkzLabelAngle[pos=0.6](C,A,D){$45^{\circ}$}
      \begin{scope}[dim style/.style={dashed,sloped,teal}]
        \tkzDrawSegment[dim={\pgfmathprintnumber 3,-6pt,}](A,B)
      \end{scope}

    \end{tikzpicture}
  \end{minipage}
  \begin{minipage}{0.7\textwidth}
    $AB=3, \; \angle ABC=90^{\circ}, \; \angle CAD=45^{\circ}, \; S_{\triangle ACD}=6$ \par
    求$BD_{max}$
  \end{minipage}
\end{exercise}
\begin{solution}
  \textbf{【解析】}几何求极值的问题，个人偏向于选取一个关键变量，可以确定其他参数，角度也好，边长也罢。这里我们设$\angle BAC=\alpha$，$\alpha$固定后，$AC$可以用$\alpha$表示，
  由于给出了面积，那么$AD$也可以用$\alpha$表示，$BD$就可以用余弦定理表示为含有$\alpha$的函数。

  $AC=\frac{3}{\cos \alpha}, \;\because S_{\triangle ACD}=6, \; \therefore \frac{1}{2}AC\times AD\sin \frac{\pi}{4}=6 \;\;\Rightarrow AD=4\sqrt{2}\cos \alpha$，根据余弦定理
  $BD^2=3^2+AD^2-2\times 3\times AD\cos(\alpha+\pi/4)=32\cos^2\alpha+9-24\sqrt{2}\cos\alpha\cos(\alpha+\pi/4)=9+8\cos^2\alpha+24\cos\alpha\sin\alpha=4\cos2\alpha+12\sin2\alpha+13=4(\cos2\alpha+3\sin2\alpha)+13$，柯西知$BD^2\le 4\times\sqrt{(1^2+3^2)}\sqrt{(\cos^22\alpha+\sin^22\alpha)}+13=13+4\sqrt{10}$，取等的时候$1:3=\cos2\alpha:\sin2\alpha$，可以取到。按惯例，$BD$根号里面还有根号，可以进一步化简为$2\sqrt{2}+\sqrt{5}$。

  关于选取变量，也可以选取$BC$长度。也有人建立直角坐标系去解。
\end{solution}

\hrulefill
\begin{exercise}
  \begin{minipage}{0.4\textwidth}
    \begin{tikzpicture}[scale=0.6]
      \tkzDefPoints{0/0/D, 0/3/A, -6/0/B, 4/0/C, -2/2/E}
      \tkzDrawSegments(A,B B,C A,C A,D C,E)
      \tkzLabelPoints(B,C,D,E)
      \tkzLabelPoints[above right](A)

      \tkzMarkAngles[scale=0.8](C,E,A)
      \tkzLabelAngle[pos=1.4, red, scale=0.8](C,E,A){$45^{\circ}$}

      \tkzMarkAngles[size=1.2, mark=|](A,C,E)
      \tkzMarkAngles[size=1.2, mark=|](E,C,B)

      \begin{scope}[dim style/.style={dashed,sloped,teal}]
        \tkzDrawSegment[dim={\pgfmathprintnumber 6,-6pt,}](B,D)
      \end{scope}
      \begin{scope}[dim style/.style={dashed,sloped,teal}]
        \tkzDrawSegment[dim={\pgfmathprintnumber 4,-6pt,}](D,C)
      \end{scope}
      \tkzMarkRightAngle[](B,D,A)
    \end{tikzpicture}
  \end{minipage}
  \begin{minipage}{0.6\textwidth}
    $BD=6, \; DC=4, \; AD \perp BC, \angle AEC=45^{\circ}, \; CE \text{平分}\angle ACB$

    求$BE$的长度。
  \end{minipage}
\end{exercise}
\begin{solution}
  \textbf{【解析】}Let $\angle ACE=\alpha$, 一旦确定$\alpha$后，$\angle B, \angle A, AD$几乎所有都能能定，所以$\alpha$应该是个关键，最后$\triangle BEC$利用正弦定理可以求出$BE$。$AD:BD=\tan(45^{\circ}-\alpha), AD:DC=\tan2\alpha \Rightarrow \tan2\alpha:\tan(45^{\circ}-\alpha)=3:2 \Rightarrow \tan\alpha=1/3, 3$，很明显$\alpha$小于$45^{\circ}$，所以$\tan \alpha=1/3$.
\end{solution}

\hrulefill
、
\begin{exercise}
  \begin{minipage}{0.3\textwidth}
    \centering
    \begin{tikzpicture}[scale=1]
      \tkzDefPoints{0/2/A, 0/0/B, 3/0/C, 3/2/D, .7/0/E}
      \tkzDefPointBy[rotation=center A angle 45](E) \tkzGetPoint{C'}
      \tkzInterLL(A,C')(C,D) \tkzGetPoint{F}
      \tkzDrawPolygons(A,B,C,D A,E,F)
      \tkzMarkAngles[size=.4, mksize=2pt,red](E,A,F)
      \tkzLabelAngle[pos=1, scale=.7, red](E,A,F){$45^{\circ}$}
      \begin{scope}[dim style/.style={dashed,sloped,red}]
        \tkzDrawSegment[dim={\pgfmathprintnumber 4,-4pt,}](A,B)
        \tkzDrawSegment[dim={\pgfmathprintnumber 6, 4pt,}](A,D)
      \end{scope}

      \tkzLabelPoints(E)
      \tkzLabelPoints[left](A,B)
      \tkzLabelPoints[right](C,D,F)
    \end{tikzpicture}
  \end{minipage}
  \begin{minipage}{0.7\textwidth}
    矩形$ABCD$，$AB=4, \; AD=6$，$E$为$BC$上一点，$\angle EAF=45^{\circ}$\\求$S_{\triangle AEF}$的最小值。
  \end{minipage}
\end{exercise}
\begin{solution}
  \textbf{【解析】}设$\angle BAE=\alpha$，那么$\angle FAD=45^{\circ}-\alpha$,$S=\frac{1}{2}AE\cdot AF\sin 45^{\circ}=\frac{1}{2}\frac{4}{\cos \alpha}\frac{6}{\cos (45^{\circ}-\alpha)}\sin 45^{\circ}$
\end{solution}

\clearpage
\section{Solutions}
\printsolutions

\clearpage
\end{document}
%%% Local Variables:
%%% mode: LaTeX
%%% TeX-master: "../main"
%%% End:
